\documentclass[12pt]{article}
\usepackage{amsmath, amssymb}

\title{Lecture 18: Marginal PMF, Correlation, Covariance, Joint Gaussian Random Variables}
\date{}

\begin{document}

\maketitle

\section{Marginal PMF}

\subsection{Definition}

For discrete random variables $X$ and $Y$, the joint PMF is:

\[
p_{XY}(x, y) = P(X = x, Y = y)
\]

\textbf{Derivation of Marginal PMFs}

The marginal PMFs are obtained by summing over the other variable:

\[
p_X(x) = \sum_{y} p_{XY}(x, y)
\]
\[
p_Y(y) = \sum_{x} p_{XY}(x, y)
\]

\textbf{Key Idea (from slide)}

The marginal PMF of one variable is obtained by summing the joint PMF over all possible values of the other variable.

\subsection{Marginal PMF from a Joint PMF Table}

From the joint PMF table:

Row sums $\rightarrow p_X(x)$ \\
Column sums $\rightarrow p_Y(y)$

Formally:

\[
p_X(x_i) = \sum_j p_{ij}, \quad p_Y(y_j) = \sum_i p_{ij}
\]

Also:

\[
\sum_i p_X(x_i) = 1, \quad \sum_j p_Y(y_j) = 1
\]

This establishes that marginal PMFs are valid probability distributions.

\subsection{Solved Example}

Given Joint PMF Table (page 5)

\textbf{Step-by-Step Solution}

Compute $p_X(x)$

\[
p_X(0) = \frac{1}{8} + \frac{3}{8} = \frac{1}{4}
\]

\[
p_X(1) = \frac{3}{8} + \frac{5}{8} = \frac{1}{4}
\]

Compute $p_Y(y)$

\[
p_Y(0) = \frac{1}{8} + \frac{2}{8} = \frac{3}{8}
\]

\[
p_Y(1) = \frac{1}{4} + \frac{1}{4} = \frac{1}{2}
\]

Final Marginal PMFs:

\[
p_X(x) =
\begin{cases}
\frac{3}{8}, & x = 0 \\
\frac{5}{8}, & x = 1
\end{cases}
\]

\[
p_Y(y) =
\begin{cases}
\frac{1}{2}, & y = 0 \\
\frac{1}{2}, & y = 1
\end{cases}
\]

\section{Correlation}

\subsection{Definition}

A basic way to measure joint behavior is through the product moment:

\[
R_{XY} = E[XY]
\]

\textbf{Interpretation}

It gives the average value of the product $XY$.

It is referred to as a raw correlation measure.

\subsection{Limitations of $R_{XY}$}

From the slides:

\begin{itemize}
\item Depends on the scale of $X$ and $Y$
\item Computed about the origin, not about the means
\end{itemize}

Therefore, it mixes:

\begin{itemize}
\item Mean values
\item Spread
\item Dependence
\end{itemize}

\textbf{Important Observation}

\[
R_{XY} = 0 \Rightarrow \text{orthogonal about the origin}
\]

\subsection{Transition to Covariance}

To remove the effect of means:

\begin{itemize}
\item Center the variables first
\end{itemize}

This leads to the concept of covariance.

\section{Covariance}

\subsection{Definition and Interpretation}

Covariance is defined as:

\[
\text{Cov}(X, Y) = E[(X - \mu_X)(Y - \mu_Y)]
\]

\textbf{Reasoning Behind Definition}

\begin{enumerate}
\item Subtract the mean from each variable
\item Measure how the centered quantities vary together
\end{enumerate}

\textbf{Interpretation (from slides)}

Covariance captures whether deviations:

\begin{itemize}
\item Occur in the same direction $\rightarrow$ Positive
\item Occur in opposite directions $\rightarrow$ Negative
\item Have no linear relation $\rightarrow$ Zero
\end{itemize}

It indicates the direction of joint linear variation.

However:

\begin{itemize}
\item Its magnitude depends on units
\item It cannot measure strength of relationship
\end{itemize}

\subsection{Algebraic Form and Derivation}

Start with:

\[
\text{Cov}(X, Y) = E[(X - \mu_X)(Y - \mu_Y)]
\]

Expand the expression:

\[
= E[XY - X\mu_Y - Y\mu_X + \mu_X \mu_Y]
\]

Substitute:

\[
\mu_X = E[X], \quad \mu_Y = E[Y]
\]

\[
\text{Cov}(X, Y) = E[XY] - \mu_Y E[X] - \mu_X E[Y] + \mu_X \mu_Y
\]

\[
\text{Cov}(X, Y) = E[XY] - \mu_X \mu_Y
\]

\[
\text{Cov}(X, Y) = E[XY] - E[X]E[Y]
\]

\subsection{Properties of Covariance}

\begin{enumerate}
\item Symmetry: $\text{Cov}(X, Y) = \text{Cov}(Y, X)$
\item Variance as Special Case: $\text{Cov}(X, X) = \text{Var}(X)$
\item Linearity with Constants: $\text{Cov}(aX + b, cY + d) = ac \cdot \text{Cov}(X, Y)$
\item Independence Implies Zero Covariance: $X \perp Y \Rightarrow \text{Cov}(X, Y) = 0$
\end{enumerate}

\textbf{Important Note}

Zero covariance means uncorrelated, but does not imply independence.

\subsection{Example 1}

\textbf{Problem}

If:

\[
Y = -2X + 3
\]

Find $\text{Cov}(X, Y)$.

\textbf{Step-by-Step Solution}

Using:

\[
\text{Cov}(X, Y) = E[XY] - E[X]E[Y]
\]

Step 1: Compute $XY$

\[
XY = X(-2X + 3) = -2X^2 + 3X
\]

Step 2: Compute expectation of $Y$

\[
E[Y] = E[-2X + 3] = -2E[X] + 3
\]

Step 3: Substitute into covariance formula

\[
\text{Cov}(X, Y) = E[XY] - E[X]E[Y]
\]

\section{Pearson’s Correlation Coefficient}

\textbf{Definition}

Measures the effect of change in one variable when the other changes.

\textbf{Properties}

\begin{itemize}
\item Measures strength and direction of linear relationship
\item Normalized version of covariance
\item Unit-free
\end{itemize}

\section{Jointly Gaussian Random Variables}

\subsection{Definition}

A pair $(X, Y)$ is jointly Gaussian if:

\begin{itemize}
\item The joint density has the bivariate Gaussian form
\item The marginals are Gaussian
\end{itemize}

\subsection{Applications (from slides)}

\begin{itemize}
\item Signal processing
\item Communication systems
\item Modeling continuous noisy measurements
\end{itemize}

\subsection{Example (from slide)}

\textbf{Scenario}

$X$: Atmospheric CO$_2$ concentration \\
$Y$: Global surface temperature anomaly

Because both vary due to many random processes, their joint behavior can often be approximated by a joint Gaussian distribution.

\subsection{Algebraic Form and Interpretation}

(Referenced in slides as continuation; algebraic form and interpretation follow from the definition of joint Gaussian density.)

\section{Final Logical Flow Summary}

\begin{enumerate}
\item Start with joint PMF
\item Derive marginal PMFs via summation
\item Measure joint behavior using $E[XY]$
\item Identify limitation $\rightarrow$ depends on means
\item Introduce covariance by centering variables
\item Derive algebraic form and properties
\item Normalize $\rightarrow$ Pearson’s coefficient
\item Extend to continuous case $\rightarrow$ Joint Gaussian RVs
\end{enumerate}

\end{document}